% !TEX root =./ECE444_Practice_main.tex
%
% L28 practice set (Null Steering Lab) -- converting computed weights into
% PHASER settings, reading a null-steered sweep, quantization limits on null
% depth, a two-channel MVDR solve by hand, and choosing between the two.

\fullwidth{\textbf{Learning Objective 3.9: } Calculate null-steering weights,
implement pattern nulls on the ADALM-PHASER, and compare manual null steering
with adaptive beamforming.
\\[4pt]\normalfont\small Show your work. A \textbf{Documentation} line at the end
is required (write \textbf{None} if you did not collaborate).}

\begin{questions}

\question \textbf{From weights to GUI settings.} A null-steering calculation for
a 4-element subarray returns the complex weights
\eq{ w_1 = 0.62 - j0.13, \quad w_2 = 0.88 + j0.05, \quad w_3 = 0.88 - j0.05,
\quad w_4 = 0.62 + j0.13 }
The Phaser GUI takes element amplitude as a percentage of the largest weight and
element phase as an offset in degrees.

\begin{parts}

	\begin{minipage}[t]{0.58\linewidth}
		\part Compute the four magnitudes and convert them to the percentages you
		would type into Element Gains.
		\begin{solution}[1.3in]
		\eq{ \vert w_1 \vert &= \sqrt{0.62^2 + 0.13^2} = 0.634 \\
		     \vert w_2 \vert &= \sqrt{0.88^2 + 0.05^2} = 0.881 }
		Elements 3 and 4 have the same magnitudes as 2 and 1. The largest weight is
		$0.881$, so the percentages are $100\vert w_n\vert/0.881$.
		\eq{ 72\%, \quad 100\%, \quad 100\%, \quad 72\% }
		\end{solution}
	\end{minipage}\hfill%
	\begin{minipage}[t]{0.37\linewidth}
		\raggedleft \vspace{12pt}
		gains = \ansbox[1.6in]{$72, 100, 100, 72\ \%$}
	\end{minipage}

	\begin{minipage}[t]{0.58\linewidth}
		\part Compute the four phase offsets in degrees.
		\begin{solution}[1.2in]
		Each phase is $\angle w_n = \arctan(\text{Im}/\text{Re})$, with both weights
		in the right half plane so no quadrant correction is needed.
		\eq{ \angle w_1 &= \arctan\lp -0.13/0.62 \rp = -11.9\mydeg \\
		     \angle w_2 &= \arctan\lp +0.05/0.88 \rp = +3.3\mydeg }
		\eq{ -11.9\mydeg, \quad +3.3\mydeg, \quad -3.3\mydeg, \quad +11.9\mydeg }
		\end{solution}
	\end{minipage}\hfill%
	\begin{minipage}[t]{0.37\linewidth}
		\raggedleft \vspace{12pt}
		phases = \ansbox[1.6in]{$-11.9, +3.3, -3.3, +11.9\mydeg$}
	\end{minipage}

	\part A student enters the four gain percentages but forgets to enter the phase
	offsets. Describe the pattern they will measure and say whether a notch appears.
		\begin{solution}[1.2in]
		With the phases left at zero the array is a uniform-phase amplitude taper of
		$72, 100, 100, 72$ percent. That is a mild taper, so the beam broadens
		slightly and the sidelobes drop by a fraction of a decibel, but the pattern
		stays symmetric about broadside and no notch appears at the intended angle.
		The amplitudes alone cannot place a null; the cancellation depends on the
		relative phases, and setting them all to zero removes it.
		\end{solution}

	\part All four weights are multiplied by $e^{j40\mydeg}$ before being entered.
	How does the measured pattern change, and why?
		\begin{solution}[1.1in]
		It does not change. A common phase factor multiplies the entire array sum by
		a constant of unit magnitude, and the measured quantity is
		$\vert F(\theta) \vert$. Only the phase \emph{differences} between elements
		affect the pattern, which is why the GUI accepts phase offsets rather than
		absolute phases.
		\end{solution}

\end{parts}

\newpage

\question \textbf{Reading a null-steered sweep.} A student sweeps the uniform
8-element array, freezes the trace, enters null-steering weights for a null at
$+22.5\mydeg$, and sweeps again. The two traces read as follows.

\begin{center}
\begin{tabular}{lcc}
\hline
 & \textbf{Uniform (frozen)} & \textbf{Null steered} \\
\hline
Peak level & $0.0$\db (reference) & $-1.8$\db \\
Level at $+22.5\mydeg$ & $-12.8\ \text{dBc}$ & $-21.6\ \text{dBc}$ \\
\hline
\end{tabular}
\end{center}

\begin{parts}

	\begin{minipage}[t]{0.58\linewidth}
		\part How much additional rejection did the null steering buy at
		$+22.5\mydeg$?
		\begin{solution}[0.9in]
		\eq{ \Delta = -12.8\db - \lp -21.6\db \rp = 8.8\db }
		The response at the interferer angle dropped by $8.8$\db relative to the
		frozen uniform reference.
		\end{solution}
	\end{minipage}\hfill%
	\begin{minipage}[t]{0.37\linewidth}
		\raggedleft \vspace{12pt}
		$\Delta$ = \ansbox[1.3in]{$8.8\db$}
	\end{minipage}

	\begin{minipage}[t]{0.58\linewidth}
		\part Express the notch depth relative to the null-steered array's own peak,
		which is the number a jammer-to-signal calculation needs.
		\begin{solution}[1.0in]
		Both levels are quoted against the uniform reference peak, so subtract the
		new peak level from the notch level.
		\eq{ -21.6\db - \lp -1.8\db \rp = -19.8\db }
		Relative to its own main lobe the array rejects the interferer by
		$19.8$\db, not $21.6$.
		\end{solution}
	\end{minipage}\hfill%
	\begin{minipage}[t]{0.37\linewidth}
		\raggedleft \vspace{12pt}
		depth = \ansbox[1.3in]{$-19.8\ \text{dBc}$}
	\end{minipage}

	\part A second student reports a notch at $-31\ \text{dBc}$ on the same bench with the
	same weights. The sweep's noise floor sits about $23$\db below the uniform peak.
	Is the report credible? Explain.
		\begin{solution}[1.2in]
		No. The sweep cannot display anything below its own noise floor, so no point
		on the trace can read more than about $23$\db below the uniform reference
		peak, which is about $21$\db below the null-steered array's own peak once the
		$1.8$\db main-lobe cost is accounted for. A
		reported $-31\ \text{dBc}$ is below the floor and indicates a reading taken from the
		wrong axis, a different reference, or a mis-set peak marker. The correct
		statement in that situation is that the notch reached the noise floor and its
		true depth is not observable with this measurement.
		\end{solution}

	\part The notch appears centered at $+19.7\mydeg$ rather than $+22.5\mydeg$.
	Give the most likely cause.
		\begin{solution}[1.1in]
		Phase quantization. The ADAR1000 rounds each commanded phase to the nearest
		$2.8125\mydeg$, so the vector the hardware applies is not the vector that was
		computed. The rounded vector is still a valid weight vector, and it places
		its null a step or two away from the requested angle. The sweep grid is the
		same $2.8125\mydeg$ step, which is why the displacement lands on a grid point.
		\end{solution}

\end{parts}

\newpage

\question \textbf{Null depth under coarse quantization.} The lab is repeated on a
board whose phase shifters have only $B = 2$ bits.

\begin{parts}

	\begin{minipage}[t]{0.58\linewidth}
		\part Find the phase LSB.
		\begin{solution}[0.9in]
		\eq{ \text{LSB} = \frac{360\mydeg}{2^B} = \frac{360\mydeg}{4} = 90\mydeg }
		\end{solution}
	\end{minipage}\hfill%
	\begin{minipage}[t]{0.37\linewidth}
		\raggedleft \vspace{12pt}
		LSB = \ansbox[1.3in]{$90\mydeg$}
	\end{minipage}

	\begin{minipage}[t]{0.58\linewidth}
		\part Use the rule of thumb that quantization sidelobes sit near $-6B$\db to
		estimate the deepest notch this board can hold.
		\begin{solution}[1.0in]
		\eq{ \text{QSLL} \approx -6B = -6(2) = -12\db }
		Quantization scatters roughly $12$\db of energy into the pattern, so a notch
		cannot be held deeper than about $12$\db below the peak. The 7-bit ADAR1000
		scatters far less: its rounded weights still permit a notch near $48$\db
		down, so the $20$ to $22$\db measured in the lab is the sweep's noise floor
		and not a quantization limit.
		\end{solution}
	\end{minipage}\hfill%
	\begin{minipage}[t]{0.37\linewidth}
		\raggedleft \vspace{12pt}
		depth $\approx$ \ansbox[1.3in]{$-12\db$}
	\end{minipage}

	\begin{minipage}[t]{0.58\linewidth}
		\part Round the four phases from Question 1(b) to the nearest multiple of the
		$90\mydeg$ LSB. What does the resulting weight vector do?
		\begin{solution}[1.1in]
		Every one of $-11.9\mydeg$, $+3.3\mydeg$, $-3.3\mydeg$ and $+11.9\mydeg$ is
		within $45\mydeg$ of zero, so all four round to $0\mydeg$.
		\eq{ 0\mydeg, \quad 0\mydeg, \quad 0\mydeg, \quad 0\mydeg }
		The phase information is gone entirely and the array is left as the amplitude
		taper of part 1(c), with no null at all.
		\end{solution}
	\end{minipage}\hfill%
	\begin{minipage}[t]{0.37\linewidth}
		\raggedleft \vspace{12pt}
		phases = \ansbox[1.4in]{$0, 0, 0, 0\mydeg$}
	\end{minipage}

	\part Explain why holding a deep null is more demanding of phase resolution
	than holding a beam peak.
		\begin{solution}[1.3in]
		At the beam peak all element contributions add in phase, and a small phase
		error reduces each contribution by $\cos$ of that error, which is a
		second-order loss. At a null the contributions are arranged to cancel
		exactly, so the residual is a difference of nearly equal quantities and any
		phase error appears in it directly, at first order. A $2.8\mydeg$ error costs
		roughly $0.01$\db at the peak and sets a floor near $-48$\db in the null,
		which is deeper than the lab sweep can display but a floor all the same.
		\end{solution}

\end{parts}

\newpage

\question \textbf{Two-channel MVDR by hand.} The PHASER's two digital channels
are steered so that a target off boresight presents the real steering vector
$s = [1, -1]^T$, while an interferer on boresight presents $u = [1, 1]^T$ with
power $P = 9$ relative to a per-channel noise power of $1$. The covariance seen
by the beamformer is $R = \sigma^2 I + P\,u u^T$.

\begin{parts}

	\begin{minipage}[t]{0.58\linewidth}
		\part Build $R$.
		\begin{solution}[1.1in]
		\eq{ u u^T = \begin{bmatrix} 1 & 1 \\ 1 & 1 \end{bmatrix}, \qquad
		     R = \begin{bmatrix} 1 & 0 \\ 0 & 1 \end{bmatrix}
		         + 9\begin{bmatrix} 1 & 1 \\ 1 & 1 \end{bmatrix}
		       = \begin{bmatrix} 10 & 9 \\ 9 & 10 \end{bmatrix} }
		\end{solution}
	\end{minipage}\hfill%
	\begin{minipage}[t]{0.37\linewidth}
		\raggedleft \vspace{12pt}
		$R$ = \ansbox[1.4in]{$\begin{bmatrix} 10 & 9 \\ 9 & 10 \end{bmatrix}$}
	\end{minipage}

	\begin{minipage}[t]{0.58\linewidth}
		\part Compute the MVDR weights $w = R^{-1}s / (s^T R^{-1} s)$.
		\begin{solution}[1.7in]
		The determinant is $100 - 81 = 19$, so
		\eq{ R^{-1} &= \frac{1}{19}\begin{bmatrix} 10 & -9 \\ -9 & 10 \end{bmatrix} \\
		     R^{-1}s &= \frac{1}{19}\begin{bmatrix} 19 \\ -19 \end{bmatrix}
		     = \begin{bmatrix} 1 \\ -1 \end{bmatrix} }
		\eq{ s^T R^{-1} s &= (1)(1) + (-1)(-1) = 2 \\
		     w &= \frac{1}{2}\begin{bmatrix} 1 \\ -1 \end{bmatrix}
		     = \begin{bmatrix} 0.5 \\ -0.5 \end{bmatrix} }
		\end{solution}
	\end{minipage}\hfill%
	\begin{minipage}[t]{0.37\linewidth}
		\raggedleft \vspace{12pt}
		$w$ = \ansbox[1.4in]{{$(0.5,\ -0.5)^T$}}
	\end{minipage}

	\begin{minipage}[t]{0.58\linewidth}
		\part Evaluate the beamformer's response to the target, $w^T s$, and to the
		interferer, $w^T u$.
		\begin{solution}[1.0in]
		\eq{ w^T s &= (0.5)(1) + (-0.5)(-1) = 1 \\
		     w^T u &= (0.5)(1) + (-0.5)(1) = 0 }
		Unit gain on the target, as the constraint requires, and a complete null on
		the interferer.
		\end{solution}
	\end{minipage}\hfill%
	\begin{minipage}[t]{0.37\linewidth}
		\raggedleft \vspace{12pt}
		$w^T s$, $w^T u$ = \ansbox[1.3in]{$1, \ 0$}
	\end{minipage}

	\part Which lab procedure did MVDR just reproduce, and what would $w$ become if
	the interferer were switched off?
		\begin{solution}[1.3in]
		The weights $[0.5, -0.5]^T$ subtract one channel from the other, which is
		procedure B: Beam 1 Phase set to $180\mydeg$, the monopulse difference beam.
		MVDR was not told about boresight interference; it found the difference beam
		because that is the quietest combination that still holds unit gain on the
		target. With the interferer switched off, $R = I$ and
		$w = s/(s^T s) = [0.5, -0.5]^T$, the ordinary matched beamformer for this
		look direction, which here happens to be the same vector.
		\end{solution}

\end{parts}

\newpage

\question \textbf{Choosing an approach.} Each situation below uses the same
8-element PHASER.

\begin{parts}

	\part A fixed ground emitter sits at a surveyed bearing of $+35\mydeg$ from the
	array face and does not move. The array must keep its beam on a target at
	$-10\mydeg$. Which approach would you use, and why?
		\begin{solution}[1.3in]
		Manual null steering. The interferer's angle is known and fixed, so the
		weights can be computed once and left in place, and the calculation gets all
		eight analog elements to work with rather than the two digital channels. It
		also costs no digital processing and no snapshot collection. The main lobe
		pays a fraction of a decibel because the null is far from the beam.
		\end{solution}

	\part An airborne jammer of unknown and changing bearing appears during a
	mission. Which approach would you use, and why?
		\begin{solution}[1.2in]
		MVDR. Manual null steering needs the interferer's angle before it can
		compute anything, and by the time an operator measures the bearing and
		re-enters eight weights the jammer has moved. MVDR estimates the covariance
		from the received data and re-solves every few snapshots, so the null follows
		the jammer without anyone knowing where it is.
		\end{solution}

	\begin{minipage}[t]{0.58\linewidth}
		\part A requirement calls for a $33$\db notch on a board whose coarse 3-bit
		phase shifters hold a null only to about $21$\db (Lesson 27). Taking roughly
		$6$\db of improvement per added bit, how many phase bits are needed?
		\begin{solution}[1.1in]
		\eq{ 33\db - 21\db = 12\db \quad \Rightarrow \quad
		     \frac{12}{6} = 2 \ \text{additional bits} }
		\eq{ B = 3 + 2 = 5 \ \text{bits}, \qquad
		     \text{LSB} = \frac{360\mydeg}{32} = 11.25\mydeg }
		\end{solution}
	\end{minipage}\hfill%
	\begin{minipage}[t]{0.37\linewidth}
		\raggedleft \vspace{12pt}
		$B$ = \ansbox[1.3in]{$5\ \text{bits}$}
	\end{minipage}

	\part Two jammers now illuminate the array from different bearings. State what
	happens when MVDR is asked to handle both on this hardware.
		\begin{solution}[1.3in]
		It cannot. Two digital channels give two degrees of freedom, one of which is
		spent on the look-direction constraint, leaving exactly one null. Facing two
		interferers the beamformer places its single null where it reduces total
		output power the most, which is generally a compromise between the two and a
		deep null on neither. Handling both would require a third digital channel, so
		the limit is the hybrid architecture rather than the algorithm.
		\end{solution}

\end{parts}

\vspace{1.5em}
\noindent\textbf{Documentation:} \rule[-2pt]{0.6\linewidth}{0.4pt}

\end{questions}
